% -*- coding: utf-8 -*-
\begin{quote}
\small \textbf{Instrucciones.} Salvo indicación en contra, puedes usar las reglas derivadas del tema~\ref{cap:dn-prop} (MT, DS, DNE, LEM, DeM) y CQ. Puedes hacerlas y comprobarlas en la aplicación \appdeduccion{}, que contiene estos ejercicios; su ayuda indica qué letras son constantes y cuáles variables.
\end{quote}

\section*{Nivel 1: cuantificador universal y silogística clásica}
\textit{En este nivel practicamos la eliminación del universal ($\forall\text{E}$, instanciando en las constantes relevantes) y la introducción del universal ($\forall\text{I}$ sobre parámetros arbitrarios que no aparecen en premisas).}

\begin{pregunta}\label{e4-1}
Demuestra el modo silogístico Celarent (ningún $M$ es $P$, y todo $S$ es $M$; por tanto, ningún $S$ es $P$):
\[
\forall x (M(x) \to \neg P(x)), \ \forall x (S(x) \to M(x)) \ \vdash \ \forall x (S(x) \to \neg P(x))
\]
\textit{Pista.} Abre una caja con hipótesis $S(c)$ para un parámetro arbitrario $c$. Instancia las dos premisas sobre $c$ mediante $\forall\text{E}$ y encadena los condicionales con $\to$E para derivar $\neg P(c)$. Cierra la subprueba con $\to$I y concluye con $\forall\text{I}$.
\end{pregunta}

\begin{pregunta}\label{e4-2}
Demuestra el modo silogístico Cesare (ningún $P$ es $M$, y todo $S$ es $M$; por tanto, ningún $S$ es $P$):
\[
\forall x (P(x) \to \neg M(x)), \ \forall x (S(x) \to M(x)) \ \vdash \ \forall x (S(x) \to \neg P(x))
\]
\textit{Pista.} Abre una caja con $S(c)$ y una subcaja con $P(c)$ para buscar contradicción con $M(c)$ y $\neg M(c)$.
\end{pregunta}

\begin{pregunta}\label{e4-3}
Demuestra la distributividad del cuantificador universal respecto a la conjunción:
\[
\forall x P(x) \land \forall x Q(x) \ \vdash \ \forall x (P(x) \land Q(x))
\]
\textit{Pista.} Extrae $\forall x P(x)$ y $\forall x Q(x)$ por $\land$E, instáncialos a la misma constante arbitraria $c$, junta por $\land$I y aplica $\forall\text{I}$.
\end{pregunta}

\begin{pregunta}\label{e4-4}
Demuestra el siguiente argumento, en el que el consecuente del condicional no contiene la variable $x$:
\[
\forall x (P(x) \to Q(a)), \ P(b) \ \vdash \ Q(a)
\]
\textit{Pista.} Como el universal afirma la implicación para cualquier individuo, basta instanciarlo mediante $\forall\text{E}$ sobre la constante $b$ de la segunda premisa y aplicar $\to$E.
\end{pregunta}

\begin{pregunta}\label{e4-5}
Demuestra la siguiente combinación de $\forall$E y $\leftrightarrow$E:
\[
\forall x (P(x) \leftrightarrow Q(x)), \ P(c) \ \vdash \ Q(c)
\]
\textit{Pista.} Aplica $\forall\text{E}$ sobre la constante $c$ y luego la eliminación del bicondicional ($\leftrightarrow$E) junto con la premisa $P(c)$.
\end{pregunta}

\section*{Nivel 2: cuantificador existencial y gestión de testigos}
\textit{En este nivel introducimos la caja de hipótesis con testigo provisional ($\exists\text{E}$). Recuerda: conviene abrir pronto las cajas de los existenciales con una constante fresca y asegurarse de que el testigo no aparezca en la conclusión.}

\begin{pregunta}\label{e4-6}
Demuestra el modo silogístico Ferio (ningún $M$ es $P$, y algún $S$ es $M$; por tanto, algún $S$ no es $P$):
\[
\forall x (M(x) \to \neg P(x)), \ \exists x (S(x) \land M(x)) \ \vdash \ \exists x (S(x) \land \neg P(x))
\]
\textit{Pista.} Abre una caja de $\exists\text{E}$ con un testigo fresco $c$ asumiendo $S(c) \land M(c)$. Descompón la conjunción con $\land\text{E}$, instancia la premisa universal sobre $c$ para obtener $\neg P(c)$, forma $S(c) \land \neg P(c)$ por $\land\text{I}$, generaliza con $\exists\text{I}$ y cierra la caja de $\exists\text{E}$.
\end{pregunta}

\begin{pregunta}\label{e4-7}
Demuestra el modo silogístico Datisi (todo $M$ es $P$, y algún $M$ es $S$; por tanto, algún $S$ es $P$):
\[
\forall x (M(x) \to P(x)), \ \exists x (M(x) \land S(x)) \ \vdash \ \exists x (S(x) \land P(x))
\]
\textit{Pista.} Abre una caja de $\exists\text{E}$ con testigo fresco $c$ ($M(c) \land S(c)$), extrae las partes, usa $\forall\text{E}$ sobre la primera premisa y $\to$E para obtener $P(c)$ y recombina como $S(c) \land P(c)$ para aplicar $\exists\text{I}$.
\end{pregunta}

\begin{pregunta}\label{e4-8}
Demuestra la distributividad del existencial respecto a la disyunción:
\[
\exists x (P(x) \lor Q(x)) \ \vdash \ \exists x P(x) \lor \exists x Q(x)
\]
\textit{Pista.} Abre la caja de testigo para el existencial ($P(c) \lor Q(c)$). Dentro de la caja, haz una prueba por casos sobre la disyunción ($\lor$E) para deducir $\exists x P(x) \lor \exists x Q(x)$.
\end{pregunta}

\begin{pregunta}\label{e4-9}
Demuestra que un universal y un existencial se combinan existencialmente:
\[
\forall x P(x), \ \exists x Q(x) \ \vdash \ \exists x (P(x) \land Q(x))
\]
\textit{Pista.} Abre la caja de $\exists\text{E}$ para la premisa existencial con testigo $c$ asumiendo $Q(c)$. Dentro de ella, aplica $\forall\text{E}$ sobre $c$ a la primera premisa, construye la conjunción por $\land\text{I}$, aplica $\exists\text{I}$ y cierra la caja.
\end{pregunta}

\begin{pregunta}\label{e4-10}
Demuestra la independencia de dos existenciales con cajas anidadas:
\[
\exists x P(x), \ \exists y Q(y) \ \vdash \ \exists x \exists y (P(x) \land Q(y))
\]
\textit{Advertencia.} Cada caja de testigo debe usar una constante distinta (por ejemplo, $c$ para la primera y $d$ para la segunda).
\end{pregunta}

\section*{Nivel 3: relaciones binarias y cuantificadores múltiples}
\textit{En este nivel trabajamos con predicados diádicos (relaciones entre individuos $R(x, y)$). Presta especial atención al orden de eliminación e introducción de los cuantificadores.}

\begin{pregunta}\label{e4-11}
Demuestra la transmisión existencial de una relación universalmente implicada:
\[
\forall x \forall y (R(x, y) \to S(x, y)), \ \exists x \forall y R(x, y) \ \vdash \ \forall y \exists x S(x, y)
\]
\textit{Pista.} Abre una subprueba de $\exists\text{E}$ para la segunda premisa con testigo fresco $c$ asumiendo $\forall y R(c, y)$. Dentro de la subprueba, toma un parámetro arbitrario $d$ para alcanzar la meta universal; instancia las premisas sobre $c$ y $d$ para derivar $S(c, d)$ por $\to$E, aplica $\exists\text{I}$ sobre $c$ para obtener $\exists x S(x, d)$ y luego $\forall\text{I}$ sobre $d$ para obtener $\forall y \exists x S(x, y)$. Por último, cierra la caja de $\exists\text{E}$.
\end{pregunta}

\begin{pregunta}\label{e4-12}
Demuestra que si una relación $R$ es simétrica, transitiva y serial (todo individuo se relaciona con al menos uno), entonces $R$ es reflexiva:
\[
\forall x \forall y (R(x, y) \to R(y, x)), \ \forall x \forall y \forall z ((R(x, y) \land R(y, z)) \to R(x, z)), \ \forall x \exists y R(x, y) \ \vdash \ \forall x R(x, x)
\]
\textit{Pista.} Para probar la meta universal $\forall x R(x, x)$, fija un parámetro arbitrario $a$. Instancia la premisa de serialidad sobre $a$ para obtener $\exists y R(a, y)$ y abre una subprueba de $\exists\text{E}$ con un testigo fresco $b$ asumiendo $R(a, b)$. Aplica la simetría sobre $a$ y $b$ para obtener $R(b, a)$, forma la conjunción $R(a, b) \land R(b, a)$ y aplica la transitividad sobre $a, b, a$ para derivar $R(a, a)$. Cierra la subprueba de $\exists\text{E}$ (el testigo $b$ no escapa) y concluye con $\forall\text{I}$ sobre $a$.
\end{pregunta}

\begin{pregunta}\label{e4-13}
Demuestra que si una relación es asimétrica, entonces es irreflexiva (nadie se relaciona consigo mismo):
\[
\forall x \forall y (R(x, y) \to \neg R(y, x)) \ \vdash \ \forall x \neg R(x, x)
\]
\textit{Pista.} Abre una caja con la hipótesis $R(c, c)$ y busca una contradicción.
\end{pregunta}

\begin{pregunta}\label{e4-14}
Demuestra el paso de existencial condicional a antecedente universal:
\[
\exists x (P(x) \to Q(a)) \ \vdash \ \forall x P(x) \to Q(a)
\]
\textit{Pista.} Abre una caja para la introducción del condicional ($\to$I) asumiendo $\forall x P(x)$. Dentro de ella, abre la caja de $\exists\text{E}$ para la premisa con un testigo fresco $c$ asumiendo $P(c) \to Q(a)$. Instancia el supuesto $\forall x P(x)$ sobre $c$ para obtener $P(c)$, deduce $Q(a)$ por $\to$E, cierra la caja de $\exists\text{E}$ (ya que $Q(a)$ no contiene a $c$) y finalmente cierra la caja de $\to$I.
\end{pregunta}

\begin{pregunta}\label{e4-15}
Demuestra la conmutatividad del doble cuantificador existencial mediante cajas anidadas:
\[
\exists x \exists y R(x, y) \ \vdash \ \exists y \exists x R(x, y)
\]
\textit{Pista.} Abre una caja de $\exists\text{E}$ con testigo $c$ para $\exists x$; dentro de ella, abre otra caja de $\exists\text{E}$ con testigo $d$ para $\exists y$. A partir de $R(c, d)$, aplica dos veces $\exists\text{I}$ para obtener $\exists y \exists x R(x, y)$ y cierra ambas cajas en orden inverso.
\end{pregunta}

\section*{Nivel 4: cuantificadores con negación, CQ y reducción al absurdo}
\textit{En este nivel combinamos cuantificadores con negación utilizando la regla de conversión de cuantificadores (CQ) y demostraciones clásicas por reducción al absurdo (IP).}

\begin{pregunta}\label{e4-16}
Demuestra mediante reglas básicas (sin usar la regla derivada CQ) que:
\[
\neg \exists x P(x) \ \vdash \ \forall x \neg P(x)
\]
\textit{Pista.} Para una constante arbitraria $c$, abre una hipótesis $P(c)$ con el fin de derivar $\neg P(c)$ por $\neg\text{I}$. Dentro de ella, aplica $\exists\text{I}$ sobre $c$ para obtener $\exists x P(x)$, lo que genera contradicción inmediata ($\bot$) con la premisa $\neg \exists x P(x)$. Cierra la subprueba con $\neg\text{I}$ y concluye con $\forall\text{I}$ sobre $c$.
\end{pregunta}

\begin{pregunta}\label{e4-17}
Demuestra el paso recíproco, también con reglas básicas:
\[
\forall x \neg P(x) \ \vdash \ \neg \exists x P(x)
\]
\textit{Pista.} Procede por introducción de la negación ($\neg\text{I}$) abriendo la hipótesis $\exists x P(x)$. Dentro de esa caja, abre una subprueba de $\exists\text{E}$ con testigo fresco $c$ asumiendo $P(c)$. Instancia la premisa universal sobre $c$ mediante $\forall\text{E}$ para obtener $\neg P(c)$ y obtén $\bot$; cierra la caja de $\exists$E (pues $\bot$ no contiene a $c$) y después la de $\neg$I.
\end{pregunta}

\begin{pregunta}\label{e4-18}
Demuestra el siguiente argumento (basta usar \emph{modus tollens} dentro de una caja de $\exists$E; no hace falta CQ):
\[
\forall x (P(x) \to Q(x)), \ \exists x \neg Q(x) \ \vdash \ \exists x \neg P(x)
\]
\end{pregunta}

\begin{pregunta}\label{e4-19}
Demuestra, mediante introducción de la negación ($\neg$I):
\[
\neg \exists x (P(x) \land Q(x)), \ P(a) \ \vdash \ \neg Q(a)
\]
\end{pregunta}

\begin{pregunta}\label{e4-20}
Demuestra el principio clásico conocido como la paradoja del bebedor (en todo bar hay alguien tal que, si él bebe, todos beben):
\[
\vdash \ \exists x (B(x) \to \forall y B(y))
\]
\textit{Pista.} Aplica la regla LEM con las hipótesis $\forall y B(y)$ y $\neg \forall y B(y)$. En la primera caja, para cualquier constante $c$, supón $B(c)$, reitera $\forall y B(y)$, cierra con $\to$I y aplica $\exists$I. En la segunda, obtén $\exists y \neg B(y)$ por CQ y abre una caja de $\exists$E con testigo $d$ y $\neg B(d)$; supón $B(d)$, obtén $\bot$ y, por X, $\forall y B(y)$; cierra con $\to$I para obtener $B(d) \to \forall y B(y)$ y aplica $\exists$I antes de cerrar la caja de $\exists$E. (Obsérvese que la validez de este principio depende de que el dominio no sea vacío: en un bar vacío no hay nadie que cumpla la condición).
\end{pregunta}

\section*{Nivel 5: análisis de falacias formales}
\textit{A continuación se presentan cuatro demostraciones aparentemente correctas, pero que contienen un error de razonamiento por violar alguna de las restricciones de $\forall$I o de $\exists$E. Para cada una de ellas:}
\begin{enumerate}
    \item Identifica la línea en la que se aplica una regla sin respetar sus restricciones.
    \item Enuncia con precisión qué restricción ha sido violada.
    \item Proporciona un contraejemplo informal en lenguaje natural que demuestre que el argumento es inválido.
\end{enumerate}

\begin{pregunta}[Falacia de generalización indebida]\label{e4-21}

\begin{center}
$\begin{nd}[rules=myrules]
    \hypo [1] {1} {P(a)} \by{Prem.}{}
    \hypo [2] {2} {\forall x (P(x) \to Q(x))} \by{Prem.}{}
    \have [3] {3} {P(a) \to Q(a)} \Ae{2}
    \have [4] {4} {Q(a)} \ce{3, 1}
    \have [5] {5} {\forall x Q(x)} \Ai{4}
\end{nd}$
\end{center}
\end{pregunta}

\begin{pregunta}[Falacia de fuga de testigo]\label{e4-22}

\begin{center}
$\begin{nd}[rules=myrules]
    \hypo [1] {1} {\exists x P(x)} \by{Prem.}{}
    \hypo [2] {2} {\forall x (P(x) \to Q(x))} \by{Prem.}{}
    \open
        \hypo [3] {3} {P(c)} \by{Hip. testigo $c$}{}
        \have [4] {4} {P(c) \to Q(c)} \Ae{2}
        \have [5] {5} {Q(c)} \ce{4, 3}
    \close
    \have [6] {6} {Q(c)} \Ee{1, 3-5}
\end{nd}$
\end{center}
\end{pregunta}

\begin{pregunta}[Falacia de elección de testigo no fresco: confusión con un individuo conocido]\label{e4-23}

\begin{center}
$\begin{nd}[rules=myrules]
    \hypo [1] {1} {P(a)} \by{Prem.}{}
    \hypo [2] {2} {\exists x Q(x)} \by{Prem.}{}
    \open
        \hypo [3] {3} {Q(a)} \by{Hip. testigo $a$}{}
        \have [4] {4} {P(a) \land Q(a)} \by{$\land$I}{1, 3}
        \have [5] {5} {\exists x (P(x) \land Q(x))} \Ei{4}
    \close
    \have [6] {6} {\exists x (P(x) \land Q(x))} \Ee{2, 3-5}
\end{nd}$
\end{center}
\end{pregunta}

\begin{pregunta}[Falacia de la inversión ilícita de cuantificadores]\label{e4-24}

\begin{center}
$\begin{nd}[rules=myrules]
    \hypo [1] {1} {\forall y \exists x R(x, y)} \by{Prem.}{}
    \have [2] {2} {\exists x R(x, d)} \Ae{1}
    \open
        \hypo [3] {3} {R(c, d)} \by{Hip. testigo $c$}{}
        \have [4] {4} {\forall y R(c, y)} \Ai{3}
        \have [5] {5} {\exists x \forall y R(x, y)} \Ei{4}
    \close
    \have [6] {6} {\exists x \forall y R(x, y)} \Ee{2, 3-5}
\end{nd}$
\end{center}
\end{pregunta}

\section*{Nivel 6: identidad}
\textit{En este nivel usamos las reglas $=$I y $=$E (sección~\ref{sec:identidad}). Recuerda que $=$E permite sustituir una, varias o todas las apariciones de un término por otro idéntico, en cualquiera de los dos sentidos.}

\begin{pregunta}[Transitividad y simetría]\label{e4-25}
$a = b, \ b = c \ \vdash \ c = a$

\textit{Pista.} Obtén primero $a = c$ con $=$E y después aplica la simetría como en el ejemplo de la simetría de la identidad (sección~\ref{sec:identidad}).
\end{pregunta}

\begin{pregunta}\label{e4-26}
$\forall x (x = a \to P(x)) \ \vdash \ P(a)$
\end{pregunta}

\begin{pregunta}[Unicidad]\label{e4-27}
Si hay exactamente un individuo con la propiedad $P$, y tanto $a$ como $b$ la tienen, entonces $a$ y $b$ son el mismo individuo:
\[
\exists x (P(x) \land \forall y (P(y) \to y = x)), \ P(a), \ P(b) \ \vdash \ a = b
\]
\textit{Pista.} Abre una caja de $\exists$E con un testigo fresco $c$, obtén $a = c$ y $b = c$, y combínalas con $=$E.
\end{pregunta}

\begin{pregunta}[Simetría como teorema]\label{e4-28}
$\vdash \forall x \forall y (x = y \to y = x)$
\end{pregunta}

%%% Local Variables:
%%% mode: latex
%%% TeX-master: "ApuntesLogica"
%%% End:
